\subsection{傅里叶反演公式}

在本节余下部分，我们假设群 G 有限，且 K 是一个代数闭域，其特征不整除 G 的阶，从而 K 的元素 \(|G|\) 不为零。

代数 K{[}G{]} 是半单的（马施克定理）且有限维。用 \(\widehat{G}\) 表示单 K{[}G{]}-模的类所成的集合。对每个 \(\lambda\in\widehat{G}\)，选取 G 的一个线性表示 \((\mathrm{V}_{\lambda},\pi_{\lambda})\)，使相伴的 K{[}G{]}-模的类为 \(\lambda\)。集合 \(\widehat{G}\) 是有限的，且向量空间 \(V_{\lambda}\) 都是有限维的（VIII, 第 141 页，例）。对任意 \(\lambda\in\widehat{G}\)，用 \(d_{\lambda}\) 表示表示 \(\pi_{\lambda}\) 的次数，即 K-向量空间 \(V_{\lambda}\) 的维数，并用 \(\chi_{\lambda}\) 表示它的特征标。

用 \(\mathrm{F}(\widehat{\mathrm{G}})\) 表示乘积代数 \(\prod_{\lambda\in\widehat{G}}\operatorname{End}_{\mathrm{K}}(\mathrm{V}_{\lambda})\)，用 \(\overline{F}\) 表示从 K{[}G{]} 到 \(\mathrm{F}(\widehat{\mathrm{G}})\) 的映射，它由 \(\overline{\mathcal{F}}(a)=(\pi_{\lambda}(a))_{\lambda\in\widehat{G}}\) 定义。由于域 K 是代数闭的，映射 \(\overline{F}\) 是一个代数同构（同前引文）。

对每个 \(\lambda \in \widehat{\mathbf{G}}\)，代数 \(\mathrm{End}_{\mathrm{K}}(\mathrm{V}_{\lambda})\) 的维数是 \(d_{\lambda}^{2}\)；代数 \(\mathrm{K}[\mathrm{G}]\) 的维数是 \(\mathrm{Card}(\mathrm{G})\)。因此我们有关系式

\[
\mathrm{Card(G)} = \sum_ {\lambda \in \widehat {\mathrm{G}}} d _ {\lambda} ^ {2}.\tag{8}
\]

用 \(\tau\) 表示代数 K{[}G{]} 中的迹；按定义，迹 \(\tau(a)\)（元素 \(a\) 的迹，此元素属于 K{[}G{]}）是自同态 \(x \mapsto ax\)（K{[}G{]} 的自同态）的迹（III, §7, No.~7, 第 515 页）。设 \(a = \sum_{g \in G} a_g g\) 是 K{[}G{]} 的一个元素；由公式 (2)，我们有 \(\tau(ag^{-1}) = |G| a_g\)（对每个 \(g \in G\)），从而有关系式

\[
a = | \mathrm{G} | ^ {- 1} \sum_ {g \in \mathrm{G}} \tau (a g ^ {- 1}) g.\tag{9}
\]

用 \(\widehat{\tau}\) 表示代数 \(F(\widehat{G})\) 中的迹。设 \(A = (A_{\lambda})_{\lambda \in \widehat{G}}\) 是 \(F(\widehat{G})\) 的一个元素；我们有（参见 III, §9, No.~3, 第 545 页，例 3）

\[
\widehat {\tau} (\mathrm{A}) = \sum_ {\lambda \in \widehat {\mathrm{G}}} d _ {\lambda} \operatorname{Tr} (\mathrm{A} _ {\lambda}).\tag{10}
\]

由于映射 \(\overline{F}\) 是一个 K-代数同构，我们有 \(\widehat{\tau}\circ\overline{F}=\tau\)，因此

\[
\tau (a) = \widehat {\tau} (\overline {{\mathcal {F}}} (a)) = \widehat {\tau} \big ((\pi_ {\lambda} (a)) _ {\lambda \in \widehat {\mathrm{G}}} \big) = \sum_ {\lambda \in \widehat {\mathrm{G}}} d _ {\lambda} \operatorname{Tr} (\pi_ {\lambda} (a)) = \sum_ {\lambda \in \widehat {\mathrm{G}}} d _ {\lambda} \operatorname{Tr} _ {\lambda} (a)
\]

对每个 \(a\in \mathrm{K}[\mathrm{G}]\) 成立，也就是

\[
\tau = \sum_ {\lambda \in \widehat {\mathrm{G}}} d _ {\lambda} \operatorname{Tr} _ {\lambda}.\tag{11}
\]

于是，由 VIII, 第 399 页的 (2)，对 \(g \in \mathbf{G}\)，我们有

\[
\sum_ {\lambda \in \widehat {G}} d _ {\lambda}   \chi_ {\lambda} (g) = \left\{ \begin{array}{l l} | G | & \text { if   } g \text {   is   the   identity   element }, \\ 0 & \text { otherwise }. \end{array} \right.\tag{12}
\]

对 \(a \in \mathrm{K}[\mathrm{G}]\)，关系式 (9) 取如下形式：

\[
a = | \mathrm{G} | ^ {- 1} \sum_ {g \in \mathrm{G}} \sum_ {\lambda \in \widehat {\mathrm{G}}} d _ {\lambda} \left. \operatorname{Tr} (\pi_ {\lambda} (a) \pi_ {\lambda} (g ^ {- 1})) g; \right.\tag{13}
\]

由此得出，对元素 \(A = (A_{\lambda})_{\lambda \in \widehat{G}}\)（\(F(\widehat{G})\) 的每个元素），我们有

\[
\overline {{\mathcal {F}}} ^ {- 1} (\mathrm{A}) = | \mathrm{G} | ^ {- 1} \sum_ {g \in \mathrm{G}} \sum_ {\lambda \in \widehat {\mathrm{G}}} d _ {\lambda} \operatorname{Tr} (\mathrm{A} _ {\lambda} \pi_ {\lambda} (g ^ {- 1})) g\tag{14}
\]

（"傅里叶反演公式"）。

对 \(\mu \in \widehat{\mathrm{G}}\)，用 \(j_{\mu}:\mathrm{End}_{\mathrm{K}}(\mathrm{V}_{\mu})\longrightarrow \prod_{\lambda \in \widehat{\mathrm{G}}} \mathrm{End}_{\mathrm{K}}(\mathrm{V}_{\lambda})\) 表示满足 \(j_{\mu}(u) = (v_{\lambda})\) 的映射，其中 \(v_{\lambda} = 0\)（若 \(\lambda \neq \mu\)），而 \(v_{\mu} = u\)。由公式 (14)，我们有

\[
\overline {{\mathcal {F}}} ^ {- 1} (j _ {\mu} (u)) = | \mathrm{G} | ^ {- 1} d _ {\mu} \sum_ {g \in \mathrm{G}} \mathrm{Tr} (u \pi_ {\mu} (g ^ {- 1})) g.\tag{15}
\]

代数 \(\mathrm{F}(\widehat{\mathrm{G}})=\prod_{\lambda\in\widehat{\mathrm{G}}} \mathrm{End}_{\mathrm{K}}(\mathrm{V}_{\lambda})\) 的中心由诸族 \((a_{\lambda}1_{\mathrm{V}_{\lambda}})_{\lambda\in\widehat{\mathrm{G}}}\) 组成，其中 \((a_{\lambda})\) 是 K 的元素的一个族。它就是 \(\overline{{\mathcal{F}}}\) 作用于代数 K{[}G{]} 的中心所得到的像。于是该中心有一组基 \((e_{\lambda})_{\lambda\in\widehat{\mathrm{G}}}\)，它由关系式

\[
\pi_ {\lambda} (e _ {\mu}) = \delta_ {\lambda \mu} 1 _ {V _ {\lambda}}\tag{16}
\]

（对 \(\lambda, \mu \in \widehat{\mathbf{G}}\)）刻画，其中 \(\delta_{\lambda \mu}\) 是克罗内克符号。由公式 (15)，对每个 \(\mu \in \widehat{\mathbf{G}}\)，我们有

\[
e _ {\mu} = | \mathrm{G} | ^ {- 1} d _ {\mu} \sum_ {g \in \mathrm{G}} \chi_ {\mu} (g ^ {- 1}) g.\tag{17}
\]

这些元素满足关系式

\[
\sum_ {\lambda \in \widehat {\mathsf {G}}} e _ {\lambda} = 1, \quad e _ {\mu} ^ {2} = e _ {\mu}, \quad \mathrm{and} \quad e _ {\mu} e _ {\nu} = 0\tag{18}
\]

对一切 \(\mu, \nu \in \widehat{\mathrm{G}}\)（满足 \(\mu \neq \nu\)）；它们是 \(\mathrm{Z}(\mathrm{K}[\mathrm{G}])\) 的不可分解幂等元（VIII, 第 147 页，命题 15）

注. —— 设 \((V, \pi)\) 是 G 的一个线性表示。由马施克定理，K{[}G{]}-模 V 是半单的。对任意 \(\lambda \in \widehat{G}\)，用 \(V^{\lambda}\) 表示 \(\lambda\) 型同型分量（即 K{[}G{]}-模 V 的同型分量）；我们有 \(V = \bigoplus_{\lambda \in \widehat{G}} V^{\lambda}\)。由 VIII, 第 147 页的命题 15 与公式 (17)，与 V 的这个分解相伴的、以 \(V^{\lambda}\) 为像的 V 的投影算子等于

\[
\pi (e _ {\lambda}) = | \mathrm{G} | ^ {- 1} d _ {\lambda} \sum_ {g \in \mathrm{G}} \chi_ {\lambda} (g ^ {- 1}) \pi (g).\tag{19}
\]

把这个公式应用于 \(\pi = \pi_{\lambda}\)，我们得到 K 的元素 \(d_{\lambda} \cdot 1\) 不为零。我们将在后面（VIII, 第 420 页，推论 2）看到 \(d_{\mu}\) 整除 G 的基数。

设 \(\lambda\) 是 \(\widehat{\mathbf{G}}\) 的一个元素。映射 \(\pi_{\lambda}\)（从 K{[}G{]} 到 \(\mathrm{End}_{\mathrm{K}}(\mathrm{V}_{\lambda})\)）诱导一个 (K{[}G{]}, K{[}G{]})-双模同构，从 \(e_{\lambda}\mathrm{K}[\mathrm{G}]\) 到 \(\mathrm{End}_{\mathrm{K}}(\mathrm{V}_{\lambda})\)。由 VIII, 第 147 页的命题 15，K{[}G{]}-子模 \(e_{\lambda}\mathrm{K}[\mathrm{G}]\) 是 K{[}G{]} 的 \(\lambda\) 型同型分量。它也是 G 的右正则表示的 \(\lambda^{\vee}\) 型同型分量（VIII, 第 143 页，命题 11），同时又是双正则表示的 \(\lambda \times \lambda^{\vee}\) 型同型分量（VIII, 第 381 页，命题 4）。

